\chapter{机器学习基础知识}

本章建立理解机器学习模型所需的统一框架。重点不在手工实现算法，而在于理解模型学习了什么、训练目标如何定义、评价指标反映什么，以及如何向 AI 提出可验证的建模要求。

机器学习（Machine Learning, ML）是人工智能的核心分支。它以历史数据为依据，估计输入与输出之间尚未被明确写出的规律，并将这些规律用于预测、分类、结构发现或序列决策。在本书的 AI 协作主线中，学习者不需要从零实现每一种算法，但必须理解模型正在学习什么、训练过程依赖哪些假设，以及生成结果应当如何验证。

本章不追求把所有算法展开，而是先建立一个统一的理解框架。对于绝大多数机器学习方法，都可以反复追问三个问题：
\begin{enumerate}
    \item \textbf{模型是什么？} 即输入 $\mathbf{x}$ 如何经过模型 $f(\mathbf{x};\theta)$ 得到预测 $\hat{y}$，其中 $\theta$ 是需要学习的参数。
    \item \textbf{目标函数是什么？} 即用什么损失函数 $L(\theta)$ 衡量预测结果与真实标签之间的差距。
    \item \textbf{参数怎么求？} 即如何通过梯度下降、反向传播、最大似然估计或其他优化算法求出合适的 $\theta$。
\end{enumerate}
后续无论学习线性回归、逻辑回归、决策树、神经网络、RNN、CNN 还是 Transformer，都应回到这三个问题上来理解。

\section{监督学习：从带标签数据中学习映射}

\textbf{监督学习（Supervised Learning）}是工业界应用最广泛的机器学习范式。它的训练数据由样本特征和真实标签组成：
\[
\mathcal{D}=\{(\mathbf{x}_i,y_i)\}_{i=1}^{N},
\]
其中 $\mathbf{x}_i$ 是第 $i$ 个样本的输入特征，$y_i$ 是对应的标准答案。模型训练的目标，是学习一个函数
\[
\hat{y}=f(\mathbf{x};\theta),
\]
使得当未来出现新的输入 $\mathbf{x}_{new}$ 时，模型能够给出合理预测。

在真实项目中，监督学习常用于预测未来需求、路段通行时间、设备故障概率或事故风险。AI 可以帮助完成模型实现和实验比较，但标签定义、数据划分、评价指标与错误成本仍须由人明确。

\subsection{回归问题：预测连续数值}

当标签 $y$ 是连续数值时，问题称为\textbf{回归（Regression）}。典型例子包括销量预测、客流预测、路段旅行时间预测、资产收益率预测等。

以交通流量预测为例，特征 $\mathbf{x}$ 可以包含历史流量、星期几、是否节假日、天气和道路属性等；标签 $y$ 是未来时段的真实流量。模型输出 $\hat{y}$ 后，可以用于拥堵预警、运力安排或运行状态分析。

最简单的回归模型是线性回归：
\[
\hat{y}_i = \mathbf{w}^{T}\mathbf{x}_i + b,
\]
其中 $\mathbf{w}$ 和 $b$ 是需要从数据中学习的参数。它清楚地体现了机器学习三要素：
\begin{itemize}
    \item \textbf{模型：} 用线性函数表达特征和标签之间的关系；
    \item \textbf{目标函数：} 用平方误差衡量预测值和真实值之间的距离；
    \item \textbf{参数求解：} 通过最小二乘法或梯度下降求出 $\mathbf{w}$ 和 $b$。
\end{itemize}

\subsection{分类问题：预测离散类别}

当标签 $y$ 是有限个类别时，问题称为\textbf{分类（Classification）}。典型例子包括产品是否次品、客户是否违约、订单是否延误、车辆是否需要维修等。

以二分类为例，标签 $y\in\{0,1\}$。模型通常不直接输出类别，而是先输出属于类别 1 的概率：
\[
\hat{p}_i=P(y_i=1|\mathbf{x}_i;\theta).
\]
然后再根据阈值 $\tau$ 做分类：
\[
\hat{y}_i =
\begin{cases}
1, & \hat{p}_i \ge \tau,\\
0, & \hat{p}_i < \tau.
\end{cases}
\]
这里有一个非常重要的工程事实：\textbf{分类模型输出概率，最终类别取决于阈值}。因此，在医疗、金融风控、交通安全等不同场景中，阈值并不一定取 0.5，而应根据漏报和误报的代价进行调整。

\section{无监督学习：从无标签数据中发现结构}

并非所有数据都有标准答案。在许多工业场景中，人工标注成本很高，或者我们本来就不知道应该分成哪些类别。此时可以使用\textbf{无监督学习（Unsupervised Learning）}，从只有特征而没有标签的数据中发现结构。

\subsection{聚类：发现数据中的群组}

聚类算法旨在把样本分为若干簇，使得同一簇内样本相似，不同簇之间差异较大。以物流选址为例，如果已知客户经纬度和订单量，但没有服务片区标签，就可以先用 K-Means 将客户聚成若干类，再把簇中心作为前置仓或配送中心候选点。

K-Means 的目标函数是簇内平方和：
\[
\min_{\{\boldsymbol{\mu}_k\},\{r_{ik}\}}
\sum_{i=1}^{N}\sum_{k=1}^{K}r_{ik}\|\mathbf{x}_i-\boldsymbol{\mu}_k\|^2,
\]
其中 $\boldsymbol{\mu}_k$ 是第 $k$ 个簇中心，$r_{ik}\in\{0,1\}$ 表示样本 $i$ 是否被分配到簇 $k$。这同样符合“模型—目标函数—参数求解”的主线：模型是假设数据围绕若干中心聚集，目标函数是最小化簇内距离，求解过程是在“分配样本”和“更新中心”之间交替迭代。

\section{训练：目标函数从哪里来}

机器学习训练本质上是一个最优化问题。给定模型 $f(\mathbf{x};\theta)$ 后，训练就是寻找参数 $\theta$，使训练集上的平均损失最小：
\[
\min_{\theta}\; L(\theta)=\frac{1}{N}\sum_{i=1}^{N}\ell(f(\mathbf{x}_i;\theta),y_i).
\]
关键问题是：不同任务为什么会选择不同的损失函数？

\subsection{回归损失：平方误差}

对于回归问题，最常见损失是均方误差（Mean Squared Error, MSE）：
\[
L_{\mathrm{MSE}}(\theta)=\frac{1}{N}\sum_{i=1}^{N}(y_i-\hat{y}_i)^2.
\]
平方误差有三个直观含义。第一，它惩罚预测值和真实值之间的距离；第二，大误差会被平方放大，因此模型会更重视明显偏离的样本；第三，在线性回归和高斯噪声假设下，最小化平方误差等价于最大化观测数据出现的概率。

以一元线性回归为例，模型为 $\hat{y}=wx+b$。不同的 $(w,b)$ 对应不同的损失值。如果把 $w$ 和 $b$ 作为横纵坐标，把损失作为高度，就可以得到一个损失曲面；在等高线图中，梯度下降就是沿着损失下降最快的方向逐步移动，直到接近最优参数。这也是许多教材用等高线解释线性回归训练过程的原因。

\subsection{二分类损失：极大似然与交叉熵}

分类问题的目标函数更容易被误解。以二分类为例，模型输出 $\hat{p}_i=P(y_i=1|\mathbf{x}_i;\theta)$。如果真实标签 $y_i=1$，我们希望 $\hat{p}_i$ 尽可能大；如果真实标签 $y_i=0$，我们希望 $1-\hat{p}_i$ 尽可能大。把两种情况合在一起，单个样本被模型“正确解释”的概率可以写为：
\[
\hat{p}_i^{y_i}(1-\hat{p}_i)^{1-y_i}.
\]
所有样本联合起来，就得到似然函数：
\[
\mathcal{L}(\theta)=\prod_{i=1}^{N}\hat{p}_i^{y_i}(1-\hat{p}_i)^{1-y_i}.
\]
训练分类模型的思想，就是让真实标签在模型下出现的概率尽可能大，即\textbf{极大似然估计（Maximum Likelihood Estimation, MLE）}。为了便于优化，通常对似然函数取负对数，得到二分类交叉熵损失：
\[
L_{\mathrm{CE}}(\theta)=
-\frac{1}{N}\sum_{i=1}^{N}
\left[
y_i\log \hat{p}_i+(1-y_i)\log(1-\hat{p}_i)
\right].
\]
这个式子的含义很直接：真实标签为 1 时，模型若给类别 1 很低概率，就会受到很大惩罚；真实标签为 0 时，模型若给类别 0 很低概率，也会受到很大惩罚。

可以通过一个简单问题来理解 MLE。假设某个样本的真实标签是 \(y=1\)，如果希望模型更准确，应该让模型预测的 \(P(y=1|\mathbf{x})\) 更大，还是让 \(P(y=0|\mathbf{x})\) 更大？答案显然是前者。反过来，如果真实标签是 \(y=0\)，就应让 \(P(y=0|\mathbf{x})=1-\hat{p}\) 更大。极大似然法说的正是这件朴素的事情：调整参数 \(\theta\)，让真实发生的标签在模型看来尽可能“有可能发生”。交叉熵只是把这个思想写成了便于优化的负对数形式。

\subsection{多分类损失：Softmax 与交叉熵}

对于 $K$ 类分类问题，模型通常输出每个类别的得分 $s_1,\dots,s_K$，再用 Softmax 转换为概率：
\[
\hat{p}_{ik}=
\frac{\exp(s_{ik})}{\sum_{j=1}^{K}\exp(s_{ij})}.
\]
若样本 $i$ 的真实类别为 $y_i$，多分类交叉熵为：
\[
L_{\mathrm{multi}}(\theta)=
-\frac{1}{N}\sum_{i=1}^{N}\log \hat{p}_{i,y_i}.
\]
也就是说，模型训练的目标是让真实类别对应的预测概率尽可能大。

这里的思想和二分类交叉熵完全一致，只是类别从两个变成了 $K$ 个。模型先为每个类别输出一个得分 $s_1,\ldots,s_K$，这些得分本身还不是概率：它们可以为正，也可以为负，并且不要求相加为 1。Softmax 的作用，就是把这些任意实数得分转换成一组合法概率 $\hat p_1,\ldots,\hat p_K$，使得每个 $\hat p_k\in(0,1)$，并且 $\sum_{k=1}^{K}\hat p_k=1$。

可以通过一个简单例子来理解。假设现在要识别一张图片属于三类中的哪一类：猫、狗、鸟。模型对这张图片输出三个得分：
\[
s_{\text{猫}}=2.0,\qquad
s_{\text{狗}}=1.0,\qquad
s_{\text{鸟}}=0.1.
\]
经过 Softmax 后，可能得到类似的预测概率：
\[
\hat p_{\text{猫}}=0.66,\qquad
\hat p_{\text{狗}}=0.24,\qquad
\hat p_{\text{鸟}}=0.10.
\]
如果这张图片的真实类别是“猫”，那么我们当然希望模型给“猫”的概率越大越好。此时交叉熵损失只关心真实类别对应的预测概率，也就是
\[
-\log \hat p_{\text{猫}}=-\log 0.66.
\]
如果模型更自信、更准确，把“猫”的概率提高到 $0.90$，那么损失变为
\[
-\log 0.90,
\]
这个值会更小。反过来，如果模型把“猫”的概率只给到 $0.05$，虽然它可能把“狗”或“鸟”的概率给得很高，但由于真实类别是“猫”，损失会变成
\[
-\log 0.05,
\]
这个值会很大。也就是说，模型真正受到惩罚的地方在于：它没有把足够大的概率分配给真实发生的类别。

因此，多分类交叉熵可以理解为：对每个样本，只取出真实类别 $y_i$ 对应的预测概率 $\hat p_{i,y_i}$，然后让这个概率尽可能大。由于优化时通常是最小化损失，所以要最小化 $-\log \hat p_{i,y_i}$。这里加上负号，是因为 $\log \hat p_{i,y_i}$ 越大表示真实类别越可能，而优化问题通常写成“最小化”的形式，所以把最大化真实类别的对数概率，等价地写成最小化负对数概率。

从极大似然的角度看，这件事也很自然。对于样本 $i$，如果真实类别是 $y_i$，那么模型认为这个标签发生的概率就是 $\hat p_{i,y_i}$。整批数据的似然就是所有真实标签概率的乘积：
\[
\prod_{i=1}^{N}\hat p_{i,y_i}.
\]
极大似然法希望这个乘积尽可能大，也就是让每一个真实标签在模型看来都尽可能“合理”。取对数后，乘积变成求和；再加上负号和平均，就得到多分类交叉熵：
\[
L_{\mathrm{multi}}(\theta)=
-\frac{1}{N}\sum_{i=1}^{N}\log \hat p_{i,y_i}.
\]
所以，多分类交叉熵并不是一个神秘的新损失函数。它本质上仍然是在做一件非常朴素的事情：如果样本真实属于第 $y_i$ 类，就调整模型参数，让第 $y_i$ 类的预测概率尽可能大。Softmax 负责把模型输出的类别得分变成概率，交叉熵负责惩罚模型没有把足够大的概率分给真实类别。

\section{参数求解：梯度下降与反向传播}

有了目标函数之后，还需要求参数。许多简单模型有解析解，例如普通线性回归可以用最小二乘公式求解。但现代机器学习模型通常参数多、结构复杂，更常使用梯度下降类算法。

\subsection{梯度下降}

梯度下降的迭代形式为：
\[
\theta_{t+1}=\theta_t-\alpha\nabla L(\theta_t),
\]
其中 $\alpha$ 是学习率，$\nabla L(\theta_t)$ 是当前参数下损失函数的梯度。直观地说，梯度指出损失上升最快的方向，因此沿着负梯度方向更新参数，就能让损失尽量下降。

在实际训练中，通常不会每次都用全部样本计算梯度，而是使用小批量随机梯度下降（Mini-batch SGD）：每次抽取一个 batch 估计梯度，然后更新参数。Adam、RMSProp、Momentum 等优化器，都是在这一基本思想上改进更新方向和步长。

% 【图示占位】源稿图示待按本书风格重绘。

相关图示 的读图重点是：等高线越靠中心表示损失越低，每一个箭头代表一次参数更新。梯度下降并不是“直接跳到最优点”，而是反复根据当前位置的梯度方向修正参数；学习率过大可能跨过低损失区域，学习率过小则会让路径变得很长。

\subsection{反向传播的统一视角}

神经网络的参数很多，且一个参数可能通过多条计算路径影响最终损失。反向传播（Backpropagation）的本质，就是利用链式法则，把损失对每个参数的影响沿计算图一层层传回去。

可以用一个统一视角理解后续模型：
\begin{itemize}
    \item \textbf{普通前馈网络：} 误差沿网络层从输出层往输入层传播；
    \item \textbf{RNN/LSTM：} 同一组参数在多个时间步重复使用，因此梯度要沿时间轴累加，这就是随时间反向传播（BPTT）；
    \item \textbf{CNN：} 卷积核在不同空间位置共享，因此同一个卷积核参数会从许多图像位置收到梯度贡献；
    \item \textbf{Self-Attention/Transformer：} Query、Key、Value 和输出投影矩阵都参与注意力计算，最终仍然通过链式法则和梯度下降学习。
\end{itemize}
因此，复杂模型并没有脱离基本原则。它们的区别主要在于计算图更复杂、路径更多、参数共享方式不同；参数估计的核心仍然是“定义损失函数，然后用梯度更新参数”。

\section{测试与评估：模型到底好不好}

训练损失低并不代表模型真正好。我们真正关心的是模型在未见过的数据上是否可靠，这称为泛化能力。为此，需要将数据分为训练集、验证集和测试集。

\begin{itemize}
    \item \textbf{训练集（Training Set）：} 用于学习模型参数；
    \item \textbf{验证集（Validation Set）：} 用于选择超参数，例如正则化强度、树深度、学习率、隐藏层大小；
    \item \textbf{测试集（Test Set）：} 用于最终评估模型，只应在模型选择完成后使用。
\end{itemize}
如果数据量较少，可以使用 K 折交叉验证（K-fold Cross-Validation）减少一次随机划分带来的偶然性。

\subsection{回归任务的评估指标}

回归任务常用指标包括：
\begin{itemize}
    \item \textbf{MAE（Mean Absolute Error）：}
    \[
    \mathrm{MAE}=\frac{1}{N}\sum_{i=1}^{N}|y_i-\hat{y}_i|.
    \]
    它表示平均绝对误差，单位和原始标签相同，解释性强。
    \item \textbf{MAPE（Mean Absolute Percentage Error）：}
    \[
    \mathrm{MAPE}=\frac{1}{N}\sum_{i=1}^{N}\left|\frac{y_i-\hat{y}_i}{y_i}\right|\times 100\%.
    \]
    它表示平均绝对百分比误差，衡量预测误差相对于真实值的比例。例如，MAPE 为 $8\%$ 通常可以理解为预测值平均偏离真实值约 $8\%$。因此，MAPE 适合需求预测、销量预测、交通流量预测等更关心相对误差的场景。但当真实值 $y_i$ 接近 0 或等于 0 时，MAPE 会变得不稳定，使用时需要特别注意。
    \item \textbf{MSE/RMSE：}
    \[
    \mathrm{MSE}=\frac{1}{N}\sum_{i=1}^{N}(y_i-\hat{y}_i)^2,\quad
    \mathrm{RMSE}=\sqrt{\mathrm{MSE}}.
    \]
    RMSE 对大误差更敏感，适合缺货、延误等大偏差代价较高的场景。
    \item \textbf{$R^2$：}
    \[
    R^2=1-\frac{\sum_i(y_i-\hat{y}_i)^2}{\sum_i(y_i-\bar{y})^2}.
    \]
    它衡量模型相对于“只预测均值”的改进程度。
\end{itemize}

\subsection{分类任务的混淆矩阵}

二分类评估通常从混淆矩阵开始：
\[
\begin{array}{c|cc}
& \text{预测为 1} & \text{预测为 0}\\ \hline
\text{真实为 1} & TP & FN\\
\text{真实为 0} & FP & TN
\end{array}
\]
其中 TP 是真正例，FP 是假正例，TN 是真负例，FN 是假负例。常用指标包括：
\[
\mathrm{Accuracy}=\frac{TP+TN}{TP+FP+TN+FN},
\]
\[
\mathrm{Precision}=\frac{TP}{TP+FP},\quad
\mathrm{Recall}=\frac{TP}{TP+FN}.
\]
在类别极度不平衡时，准确率可能具有误导性。例如 1\% 的样本是故障，如果模型永远预测“无故障”，准确率也有 99\%，但它完全没有业务价值。

% 【图示占位】源稿图示待按本书风格重绘。

相关图示 把分类结果拆成四种情况。对于安全检测、疾病筛查或交通事件识别等任务，FN 和 FP 的代价往往并不相同，因此只看 Accuracy 容易掩盖业务风险。Precision 更关注“报出来的正例有多少是真的”，Recall 更关注“真实正例中有多少被找出来”。

\subsection{ROC 曲线是怎么画出来的}

许多读者对 ROC 曲线的来源不清楚。关键在于：分类模型通常先输出概率，而不是直接输出类别。给定不同阈值 $\tau$，同一个模型会产生不同的预测类别，从而得到不同的混淆矩阵。

对于每个阈值，可以计算：
\[
\mathrm{TPR}=\frac{TP}{TP+FN},
\]
\[
\mathrm{FPR}=\frac{FP}{FP+TN}.
\]
其中 TPR 也称为召回率，表示真实正样本中被识别出来的比例；FPR 表示真实负样本中被错误报为正样本的比例。把阈值从高到低不断变化，就会得到一系列 $(\mathrm{FPR},\mathrm{TPR})$ 点，将这些点连起来就是 ROC 曲线。

因此，ROC 曲线不是某一个固定混淆矩阵画出来的，而是\textbf{同一个概率模型在不同分类阈值下产生的一系列混淆矩阵}汇总出来的。AUC 则是 ROC 曲线下的面积，反映模型把正样本排在负样本前面的整体能力。

% 【图示占位】源稿图示待按本书风格重绘。

相关图示 左侧强调阈值变化带来的分类结果变化，右侧则把这些结果映射为 ROC 空间中的点。阈值较高时，模型更保守，通常 FPR 较低但 TPR 也可能较低；阈值降低后，更多样本会被判为正类，TPR 和 FPR 都可能上升。

\section{欠拟合、过拟合与解决方法}

模型训练中最常见的两个问题是欠拟合和过拟合。

% 【图示占位】源稿图示待按本书风格重绘。

相关图示 展示的是模型复杂度与泛化能力之间的折中。欠拟合时，模型连主要趋势都没有学到；合适拟合时，模型捕捉主趋势但不过度追逐噪声；过拟合时，训练集上的点可能被贴得很紧，但对新样本的预测反而不稳定。

\subsection{欠拟合}

\textbf{欠拟合（Underfitting）}指模型过于简单，连训练数据中的主要规律都没有学到。例如用直线拟合明显的非线性需求曲线，或者用太浅的树模型处理复杂分类边界。欠拟合通常表现为训练误差和测试误差都很高。

常见解决方法包括：
\begin{itemize}
    \item 增加模型复杂度，例如加入非线性特征、提高多项式阶数、增加树深度或网络宽度；
    \item 增加有解释力的特征；
    \item 降低过强的正则化；
    \item 训练更久，或调整学习率使优化更充分。
\end{itemize}

\subsection{过拟合}

\textbf{过拟合（Overfitting）}指模型把训练数据中的噪声也记住了，导致训练集表现很好，但测试集表现变差。过拟合在 AI 协作建模中尤其危险：如果预测模型把随机波动当成真实规律，使用者可能在缺少外部验证的情况下接受并传播错误结论。

常见解决方法包括：
\begin{itemize}
    \item 增加训练数据，或使用数据增强；
    \item 使用正则化，例如岭回归的 $L_2$、Lasso 的 $L_1$、神经网络中的 weight decay；
    \item 降低模型复杂度，例如限制树深度、减少神经网络层数；
    \item 使用交叉验证选择超参数；
    \item 使用 Early Stopping，在验证集表现开始变差前停止训练；
    \item 使用集成方法，例如随机森林通过 bagging 降低单棵树的方差。
\end{itemize}

\section{强化学习简介}

监督学习和无监督学习主要从静态数据中学习规律，而\textbf{强化学习（Reinforcement Learning, RL）}关注智能体如何在环境中通过试错学习策略。它的基本元素包括状态 $s$、动作 $a$、奖励 $r$、策略 $\pi(a|s)$ 和环境转移。

强化学习与真实工程决策中的马尔可夫决策过程（MDP）和动态规划有密切关系。例如，仓储机器人在仓库中移动，每一步消耗电量，撞到障碍物受到惩罚，到达目标点获得奖励。智能体通过反复试错，学习在不同状态下应该采取什么动作。

强化学习会在后续章节中展开。本章只需把它理解为机器学习的第三类范式：监督学习学习“输入到标签”的映射，无监督学习发现数据结构，强化学习学习“状态到动作”的策略。

\section{案例占位}

\begin{quote}
\textbf{【案例占位】}

后续将在此加入一个围绕。案例不以逐行教学代码为主，而将展示问题定义、向 AI 提供的上下文与约束、模型选择依据、关键中间产物、验证过程以及人需要承担的判断。

\textbf{【待确定内容】} 数据来源、任务背景、AI 协作流程、模型比较、错误诊断与现实解释。
\end{quote}
